%=========================================================
% Notas para el curso de Calculo Integral - TNM
% Formato conservado de Unidad2_Calculo_Integral_TNM
%=========================================================
\documentclass[12pt]{article}
\usepackage[utf8]{inputenc}
\usepackage[T1]{fontenc}
\usepackage[spanish,mexico]{babel}
\usepackage{lmodern}
\usepackage{amsmath,amssymb,amsthm,mathtools}
\usepackage{tikz}
\usepackage{pgfplots}
\pgfplotsset{compat=1.18}
\usepgfplotslibrary{fillbetween}
\usepackage[most]{tcolorbox}
\usepackage{multicol}
\usepackage{needspace}
\usepackage{geometry,enumitem,booktabs,array,xcolor,graphicx,hyperref,fancyhdr,microtype,longtable}
\geometry{letterpaper,top=2.3cm,bottom=2.3cm,left=2.5cm,right=2.5cm}
\setlength{\parindent}{0pt}
\setlength{\parskip}{0.55em}
\setlength{\headheight}{15pt}
\setlength{\emergencystretch}{2em}
\definecolor{azul}{RGB}{25,70,120}
\hypersetup{colorlinks=true,linkcolor=azul,urlcolor=azul,citecolor=azul,pdftitle={Calculo Integral - Unidad 4: Series},pdfauthor={Carlos Ernesto Martinez Rodriguez}}
\newtheorem{teorema}{Teorema}[section]
\newtheorem{proposicion}[teorema]{Proposición}
\newtheorem{corolario}[teorema]{Corolario}
\theoremstyle{definition}
\newtheorem{definicion}[teorema]{Definición}
\newtheorem{ejemplo}[teorema]{Ejemplo}
\newtheorem{Ejer}[teorema]{Ejercicio}
\newtheorem{nota}[teorema]{Nota}
\theoremstyle{remark}
\newtheorem*{observacion}{Observación}
\newtheorem*{advertencia}{Advertencia}
\newtheorem*{estrategia}{Estrategia}
\newtcolorbox{notateorica}[1][]{enhanced,breakable,colback=gray!5,colframe=black!65,fonttitle=\bfseries,title=#1,boxrule=0.7pt,arc=2mm,left=2mm,right=2mm,top=2mm,bottom=2mm}
\newtcolorbox{formula}[1][]{enhanced,breakable,colback=blue!3,colframe=blue!55!black,fonttitle=\bfseries,title=#1,boxrule=0.8pt,arc=2mm,left=2mm,right=2mm,top=2mm,bottom=2mm}
\newcommand{\R}{\mathbb{R}}
\newcommand{\dd}{\,dx}
\newcommand{\dt}{\,dt}
\newcommand{\du}{\,du}
\newcommand{\diff}{\mathop{}\!\mathrm{d}}
\newcommand{\F}[1]{\[\displaystyle #1\]}
\pagestyle{fancy}
\fancyhf{}
\fancyhead[L]{\small Tecnológico Nacional de México}
\fancyhead[R]{\small Cálculo Integral}
\fancyfoot[L]{\footnotesize Carlos Ernesto Martínez Rodríguez}
\fancyfoot[C]{\footnotesize\thepage}
\fancyfoot[R]{\footnotesize mrdgzcarlose@gmail.com}
\renewcommand{\headrulewidth}{0.4pt}
\renewcommand{\footrulewidth}{0.4pt}
\title{\textbf{Notas para el curso de Cálculo Integral}\\[0.3cm]\Large Unidad 4: Series}
\author{Carlos Ernesto Martínez Rodríguez\\Tecnológico Nacional de México}
\date{Agosto--Diciembre 2026}
\begin{document}
\maketitle
\tableofcontents
\newpage
%=========================================================
\section*{Presentación}
%=========================================================
En la Unidad 3 estudiamos magnitudes que pueden construirse como límites de sumas. En esta unidad invertimos la perspectiva: primero estudiaremos el comportamiento de \emph{sucesiones} y de \emph{sumas parciales}; después utilizaremos series de potencias para representar funciones y aproximar cantidades que no siempre pueden expresarse mediante operaciones elementales.

La metodología es progresiva, en el estilo de Stewart: motivar mediante un ejemplo, definir con precisión, estudiar un criterio de convergencia, observar sus hipótesis y practicar con ejercicios de estructura similar antes de combinar técnicas. Las series y los polinomios de Taylor sólo sirven como aproximaciones fiables si entendemos sus condiciones de convergencia y el error restante.

\textbf{Alcance del temario:} sucesiones; series finitas e infinitas; convergencia y divergencia; pruebas de convergencia; series de potencias y su intervalo/radio; derivación e integración; Taylor y Maclaurin; aproximación de funciones e integrales; visualización con software. Planeación: 20 horas (12 teóricas y 8 prácticas).

%=========================================================
\section{Sucesiones y su límite}
%=========================================================
Una sucesión es una función cuyo dominio son los números naturales (o los enteros a partir de algún índice $n_0$). Escribimos $a_n$ para designar su término de índice $n$ y $\{a_n\}$ para designar toda la sucesión.
\begin{definicion}[Convergencia de una sucesión]
Decimos que $\{a_n\}$ converge a $L\in\R$ cuando para todo $\varepsilon>0$ existe $N\in\mathbb N$ tal que
\[\boxed{n\ge N\implies |a_n-L|<\varepsilon.}\]
Escribimos $\lim_{n\to\infty}a_n=L$. Si no existe ningún límite real, la sucesión diverge.
\end{definicion}
Esta definición enfatiza que, a partir de cierto índice, \emph{todos} los términos permanecen arbitrariamente cerca de $L$. No basta observar que algunos términos se aproximan.
\begin{ejemplo}[Un límite de sucesión desde la definición]
Sea $a_n=(2n+1)/n=2+1/n$. Queremos probar que $a_n\to2$. Dado $\varepsilon>0$, basta exigir
\[|a_n-2|=\frac1n<\varepsilon.\]
Podemos elegir cualquier entero $N>1/\varepsilon$. Si $n\ge N$, se verifica la desigualdad requerida, luego $\boxed{\lim a_n=2}$.
\end{ejemplo}
\begin{ejemplo}[Una sucesión oscilante]
La sucesión $a_n=(-1)^n$ toma alternadamente los valores $-1$ y $1$. No converge: los términos pares permanecen en $1$ y los impares en $-1$. En cambio, $b_n=(-1)^n/n$ satisface $|b_n|=1/n\to0$; por el teorema de compresión, $b_n\to0$.
\end{ejemplo}
\begin{proposicion}[Álgebra y orden de límites]
Si $a_n\to A$, $b_n\to B$ y $c\in\R$, entonces $a_n\pm b_n\to A\pm B$, $ca_n\to cA$, $a_nb_n\to AB$ y, si $B\ne0$, $a_n/b_n\to A/B$ (para índices suficientemente grandes). Si $a_n\le b_n\le c_n$ y los extremos convergen al mismo límite $L$, entonces $b_n\to L$.
\end{proposicion}
\begin{ejemplo}[Dividir por la mayor potencia]
Para $a_n=(3n^2-n+1)/(2n^2+5)$, dividimos numerador y denominador entre $n^2$:
\[a_n=\frac{3-1/n+1/n^2}{2+5/n^2}
\longrightarrow\boxed{\frac32}.\]
Esta misma técnica aparece en límites al infinito de funciones racionales.
\end{ejemplo}
\begin{teorema}[Sucesión monótona acotada]
Toda sucesión creciente y acotada superiormente converge. Toda sucesión decreciente y acotada inferiormente converge.
\end{teorema}
\begin{proof}
Si $a_n$ es creciente y está acotada superiormente, el conjunto $\{a_n:n\in\mathbb N\}$ tiene supremo $L$ por completitud de $\R$. Para cada $\varepsilon>0$, existe algún $N$ con $a_N>L-\varepsilon$, pues si no $L-\varepsilon$ sería cota superior. Por monotonía, $L-\varepsilon<a_N\le a_n\le L$ para $n\ge N$, así que $|a_n-L|<\varepsilon$. El argumento decreciente es análogo, utilizando el ínfimo.
\end{proof}
\begin{ejemplo}[Una sucesión definida recursivamente]
Defina $a_1=1$ y $a_{n+1}=\sqrt{2+a_n}$. Si $a_n\le2$, entonces $a_{n+1}\le2$; por inducción todos los términos están acotados superiormente por $2$. Además, para $1\le a_n\le2$ se cumple $a_{n+1}\ge a_n$, pues $2+a_n\ge a_n^2$ en dicho intervalo. Así la sucesión es creciente y acotada, por lo que converge a un límite $L\in[1,2]$. Pasando al límite en la recurrencia,
\[L=\sqrt{2+L}\implies L^2-L-2=0\implies L=2\text{ o }L=-1.\]
La condición $L\in[1,2]$ elimina $-1$, de modo que $\boxed{L=2}$.
\end{ejemplo}

%=========================================================
\section{Series finitas e infinitas}
%=========================================================
Una serie infinita no es una suma realizada instantáneamente: se define como el límite de una sucesión de \emph{sumas parciales}.
\begin{definicion}[Serie y suma parcial]
Dada una sucesión $\{a_n\}_{n\ge1}$, su suma parcial $N$-ésima es
\[S_N=\sum_{n=1}^N a_n=a_1+\cdots+a_N.\]
La serie $\sum_{n=1}^{\infty}a_n$ converge a $S$ cuando $S_N\to S$; en ese caso escribimos
\[\boxed{\sum_{n=1}^{\infty}a_n=\lim_{N\to\infty}S_N=S.}\]
Si las sumas parciales no convergen a un número real, la serie diverge.
\end{definicion}
\begin{teorema}[Condición necesaria de convergencia]
Si $\sum_{n=1}^{\infty}a_n$ converge, entonces $a_n\to0$.
\end{teorema}
\begin{proof}
Sea $S_N=\sum_{n=1}^N a_n$ y supongamos $S_N\to S$. Como $a_n=S_n-S_{n-1}$, al tomar el límite obtenemos $a_n\to S-S=0$.
\end{proof}
\begin{advertencia}
El recíproco es falso: $a_n\to0$ \emph{no garantiza} que $\sum a_n$ converja. Éste es uno de los errores más frecuentes al estudiar series.
\end{advertencia}

\subsection*{Serie geométrica}
Para $r\ne1$,
\[1+r+\cdots+r^N=\frac{1-r^{N+1}}{1-r}.\]
Esta igualdad se demuestra multiplicando la suma por $r$ y restando: casi todos los términos se cancelan. Al tomar el límite,
\begin{teorema}[Serie geométrica]
\[\boxed{\sum_{n=0}^{\infty}ar^n=\frac{a}{1-r}\qquad\text{si }|r|<1.}\]
Si $a\ne0$ y $|r|\ge1$, los términos no tienden a cero (o las sumas parciales no convergen), por lo que diverge.
\end{teorema}
\begin{ejemplo}[Una suma geométrica]
\[\sum_{n=0}^\infty\frac{3}{2^n}
=3\sum_{n=0}^\infty\left(\frac12\right)^n
=\frac3{1-1/2}=\boxed6.\]
El primer término corresponde a $n=0$, no a $n=1$; distinguir el índice inicial evita factores faltantes.
\end{ejemplo}
\begin{ejemplo}[Una serie geométrica alternante]
\[\sum_{n=1}^\infty\frac{(-1)^{n-1}}{3^n}
=\frac13\sum_{k=0}^\infty\left(-\frac13\right)^k
=\frac{1/3}{1+1/3}=\boxed{\frac14}.\]
Antes de aplicar la fórmula identificamos primer término y razón.
\end{ejemplo}

\subsection*{Series telescópicas}
\begin{ejemplo}[Cancelación de términos]
Utilizando fracciones parciales,
\[\frac1{n(n+1)}=\frac1n-\frac1{n+1}.\]
Entonces
\begin{align*}
S_N&=\sum_{n=1}^N\left(\frac1n-\frac1{n+1}\right)\\
&=\left(1-\frac12\right)+\left(\frac12-\frac13\right)+\cdots+\left(\frac1N-\frac1{N+1}\right)\\
&=1-\frac1{N+1}\longrightarrow\boxed1.
\end{align*}
\end{ejemplo}

\subsection*{La serie armónica}
Aunque $1/n\to0$, la serie $\sum_{n=1}^{\infty}1/n$ diverge. Agrupando términos,
\[1+\frac12+\left(\frac13+\frac14\right)
+\left(\frac15+\frac16+\frac17+\frac18\right)+\cdots,\]
cada nuevo grupo, con el doble de términos, suma al menos $1/2$. Hay infinitos grupos, de modo que las sumas parciales crecen sin cota. Por eso $\sum1/n$ diverge.

%=========================================================
\section{Criterios de convergencia para series positivas}
%=========================================================
Ningún criterio resuelve todos los problemas; se elige examinando la estructura del término general. Para series con $a_n\ge0$, las sumas parciales son crecientes: la convergencia equivale a que estén acotadas.

\subsection*{Prueba de la integral y series $p$}
\begin{teorema}[Criterio integral]
Sea $f:[1,\infty)\to(0,\infty)$ continua y decreciente, y sea $a_n=f(n)$. Entonces
\[\sum_{n=1}^\infty a_n\quad\text{y}\quad\int_1^\infty f(x)\dd\]
convergen o divergen simultáneamente.
\end{teorema}
\begin{proof}
La monotonía permite encerrar las áreas de los rectángulos: para $n\ge1$,
\[\int_1^{N+1}f(x)\dd\le\sum_{n=1}^N f(n)
\le f(1)+\int_1^N f(x)\dd.\]
Si la integral converge, las sumas parciales positivas están acotadas superiormente y por tanto convergen. Si la integral diverge, la cota inferior obliga a que también diverjan las sumas parciales.
\end{proof}
\begin{teorema}[Serie $p$]
\[\boxed{\sum_{n=1}^\infty\frac1{n^p}\text{ converge si y sólo si }p>1.}\]
\end{teorema}
\begin{proof}
Para $p>0$, $f(x)=x^{-p}$ es positiva, continua y decreciente; la integral $\int_1^\infty x^{-p}\dd$ converge exactamente cuando $p>1$. Si $p\le0$, los términos $n^{-p}$ no tienden a cero y la serie diverge también. Por tanto la equivalencia es válida para todos los valores reales de $p$.
\end{proof}
\begin{ejemplo}[Una serie $p$]
La serie $\sum_{n=1}^\infty n^{-3/2}$ converge, pues $3/2>1$. No necesitamos conocer su suma exacta para decidirlo. En cambio, $\sum1/\sqrt n$ diverge porque $p=1/2\le1$.
\end{ejemplo}

\subsection*{Comparación directa}
\begin{teorema}[Comparación]
Si $0\le a_n\le b_n$ para índices suficientemente grandes, entonces:
\begin{enumerate}[label=(\roman*)]
\item Si $\sum b_n$ converge, también $\sum a_n$ converge.
\item Si $\sum a_n$ diverge, también $\sum b_n$ diverge.
\end{enumerate}
\end{teorema}
\begin{ejemplo}[Comparación con una serie $p$]
Como $n^2+1\ge n^2$ para $n\ge1$,
\[0<\frac{1}{n^2+1}\le\frac1{n^2}.\]
La serie $\sum1/n^2$ converge, luego $\sum1/(n^2+1)$ también converge.
\end{ejemplo}
\begin{ejemplo}[Probar divergencia]
Para $n\ge1$, se verifica $n+1\le2n$, y por ello
\[\frac1{n+1}\ge\frac1{2n}.\]
Como $\sum1/(2n)$ diverge, se deduce que $\sum1/(n+1)$ diverge.
\end{ejemplo}

\subsection*{Comparación por el límite}
\begin{teorema}[Comparación por el límite]
Supongamos $a_n>0$, $b_n>0$ y
\[\lim_{n\to\infty}\frac{a_n}{b_n}=L,\qquad0<L<\infty.\]
Entonces $\sum a_n$ y $\sum b_n$ tienen el mismo carácter de convergencia.
\end{teorema}
\begin{ejemplo}[Reconocer el comportamiento dominante]
Estudiemos
\[\sum_{n=1}^\infty\frac{3n+1}{n^3+2}.\]
Para $n$ grande, el término es comparable con $3/n^2$. Tomamos $b_n=1/n^2$:
\[\lim_{n\to\infty}\frac{(3n+1)/(n^3+2)}{1/n^2}
=\lim_{n\to\infty}\frac{3n^3+n^2}{n^3+2}=3.\]
Puesto que el límite es positivo y finito y $\sum1/n^2$ converge, la serie original también converge.
\end{ejemplo}
\begin{observacion}
Los criterios de comparación exigen términos no negativos (al menos a partir de un índice). Si los signos alternan, puede ser útil estudiar primero la serie de valores absolutos.
\end{observacion}

%=========================================================
\section{Series alternantes, convergencia absoluta y pruebas de razón y raíz}
%=========================================================
Una serie alternante cambia de signo; por ejemplo, $1-1/2+1/3-1/4+\cdots$. La cancelación puede producir convergencia incluso cuando los valores absolutos no son sumables.
\begin{teorema}[Criterio de Leibniz]
Si $b_n\ge0$, $b_{n+1}\le b_n$ para todos los índices suficientemente grandes y $b_n\to0$, entonces
\[\sum_{n=1}^\infty(-1)^{n-1}b_n\]
converge. Si $S_N$ es una suma parcial, el error satisface
\[\boxed{|S-S_N|\le b_{N+1}}\]
y su signo coincide con el del primer término omitido.
\end{teorema}
\begin{proof}
Las sumas parciales pares cumplen $S_{2n+2}-S_{2n}=b_{2n+1}-b_{2n+2}\ge0$, por lo que forman una sucesión creciente. Las impares cumplen $S_{2n+3}-S_{2n+1}=-b_{2n+2}+b_{2n+3}\le0$; son decrecientes. Además $S_{2n}\le S_{2n+1}$ y su diferencia es $b_{2n+1}\to0$. Ambas subsecuencias convergen al mismo límite $S$. Como $S$ queda entre $S_N$ y $S_{N+1}$, su distancia a $S_N$ no supera $b_{N+1}$.
\end{proof}
\begin{ejemplo}[Convergencia de la armónica alternante]
La serie
\[1-\frac12+\frac13-\frac14+\cdots\]
converge por Leibniz, pues $b_n=1/n$ decrece a cero. Sin embargo, la serie de valores absolutos es la armónica, que diverge. Es un ejemplo de \emph{convergencia condicional}.
\end{ejemplo}
\begin{definicion}
Una serie $\sum a_n$ converge \emph{absolutamente} si $\sum|a_n|$ converge. Converge \emph{condicionalmente} si $\sum a_n$ converge, pero $\sum|a_n|$ diverge.
\end{definicion}
\begin{teorema}
La convergencia absoluta implica convergencia de la serie original.
\end{teorema}
\begin{proof}
Para las sumas finitas se tiene $|\sum_{k=n}^m a_k|\le\sum_{k=n}^m|a_k|$. Si $\sum|a_n|$ converge, sus colas pueden hacerse arbitrariamente pequeñas. La desigualdad muestra que también las colas de $\sum a_n$ son pequeñas; por el criterio de Cauchy, $\sum a_n$ converge.
\end{proof}

\subsection*{Criterio de la razón}
\begin{teorema}[Razón de d'Alembert]
Para términos no nulos, si existe
\[L=\lim_{n\to\infty}\left|\frac{a_{n+1}}{a_n}\right|,\]
entonces $\sum a_n$ converge absolutamente si $L<1$, diverge si $L>1$ (o el límite es infinito) y el criterio \emph{no decide} cuando $L=1$.
\end{teorema}
\begin{ejemplo}[Factoriales]
Para $\sum_{n=1}^{\infty}n/2^n$, tomamos $a_n=n/2^n$. Así,
\[\left|\frac{a_{n+1}}{a_n}\right|
=\frac{n+1}{2^{n+1}}\frac{2^n}{n}
=\frac{n+1}{2n}\longrightarrow\boxed{\frac12}.\]
Como $L<1$, la serie converge.
\end{ejemplo}
\begin{ejemplo}[Un criterio inconcluso]
Para $\sum1/n^2$,
\[\left|\frac{a_{n+1}}{a_n}\right|=\frac{n^2}{(n+1)^2}\longrightarrow1.\]
La serie \emph{sí converge} por el criterio $p$, pero la prueba de la razón, cuando da $1$, no demuestra ni convergencia ni divergencia.
\end{ejemplo}

\subsection*{Criterio de la raíz}
\begin{teorema}[Raíz de Cauchy]
Sea $L=\lim_{n\to\infty}\sqrt[n]{|a_n|}$, si existe. Entonces $\sum a_n$ converge absolutamente si $L<1$, diverge si $L>1$ (o $L=\infty$), y el criterio no decide cuando $L=1$.
\end{teorema}
\begin{ejemplo}[Potencias de expresiones racionales]
Estudie
\[\sum_{n=1}^\infty\left(\frac{2n}{3n+1}\right)^n.\]
La raíz $n$-ésima del término es $2n/(3n+1)\to2/3<1$. La serie converge, y no fue necesario desarrollar las potencias.
\end{ejemplo}
\begin{estrategia}
Usualmente conviene probar la \emph{razón} si aparecen factoriales o productos de muchos factores sucesivos, y la \emph{raíz} cuando todo el término está elevado a la potencia $n$. Si se obtiene $1$, cambie de criterio; no concluya que la serie diverge.
\end{estrategia}

%=========================================================
\section{Series de potencias: centro, radio e intervalo}
%=========================================================
Una serie de potencias permite representar funciones utilizando polinomios de grado creciente alrededor de un centro $c$.
\begin{definicion}
Una serie de potencias centrada en $c$ tiene la forma
\[\boxed{\sum_{n=0}^\infty a_n(x-c)^n.}\]
Su \emph{radio de convergencia} $R$ puede ser $0$, un número positivo finito o $\infty$. Si $0<R<\infty$, la serie converge absolutamente cuando $|x-c|<R$ y diverge cuando $|x-c|>R$. En los extremos $x=c\pm R$ debe analizarse por separado.
\end{definicion}
\begin{ejemplo}[Radio e intervalo de convergencia]
Determine el intervalo de convergencia de
\[\sum_{n=1}^\infty\frac{(x-2)^n}{n3^n}.\]
Para un $x$ fijo, el término es $a_n=(x-2)^n/(n3^n)$. Por razón,
\[\left|\frac{a_{n+1}}{a_n}\right|
=\frac{|x-2|}{3}\frac{n}{n+1}
\longrightarrow\frac{|x-2|}{3}.\]
Hay convergencia absoluta si $|x-2|<3$, equivalente a $-1<x<5$. Esto identifica el radio $R=3$, pero \emph{no} decide los extremos.

En $x=5$ se obtiene $\sum_{n=1}^\infty1/n$, divergente. En $x=-1$ se obtiene $\sum_{n=1}^\infty(-1)^n/n$, convergente por Leibniz, aunque no absolutamente. En consecuencia,
\[\boxed{R=3,\qquad I=[-1,5).}\]
\end{ejemplo}
\begin{ejemplo}[Extremos de una serie geométrica]
Para
\[\sum_{n=0}^\infty\frac{(x+1)^n}{4^n}\]
la razón geométrica es $r=(x+1)/4$. Se requiere $|x+1|<4$, es decir $-5<x<3$. En los extremos, los términos son $(-1)^n$ o $1$, respectivamente, y \emph{no} tienden a cero. Por tanto,
\[\boxed{R=4,\qquad I=(-5,3).}\]
\end{ejemplo}
\begin{observacion}
El radio de convergencia determina el intervalo \emph{abierto} central. La inclusión o exclusión de los extremos depende de la serie concreta. Por eso la respuesta ``$R=3$'' no sustituye la petición de encontrar el intervalo.
\end{observacion}

%=========================================================
\section{Derivación e integración de series de potencias}
%=========================================================
Los polinomios se derivan e integran término a término. Para series de potencias esto también es válido \emph{dentro del radio de convergencia}, y las series resultantes conservan el mismo radio. En los extremos se requiere un estudio independiente.
\begin{teorema}[Operaciones término a término]
Si $\sum_{n=0}^\infty a_n(x-c)^n$ tiene radio $R>0$ y define $f(x)$ para $|x-c|<R$, entonces
\[\boxed{f'(x)=\sum_{n=1}^\infty n a_n(x-c)^{n-1}}\]
y
\[\boxed{\int f(x)\dd=C+\sum_{n=0}^\infty\frac{a_n}{n+1}(x-c)^{n+1}.}\]
Ambas series poseen radio $R$.
\end{teorema}
\begin{ejemplo}[De la geométrica al logaritmo]
Para $|x|<1$ conocemos la identidad
\[\frac1{1+x}=\sum_{n=0}^\infty(-1)^n x^n.\]
Integrando entre $0$ y $x$,
\begin{align*}
\ln(1+x)&=\int_0^x\frac1{1+t}\,dt
=\sum_{n=0}^\infty(-1)^n\int_0^x t^n\,dt\\
&=\sum_{n=0}^\infty\frac{(-1)^n x^{n+1}}{n+1}
=\boxed{x-\frac{x^2}{2}+\frac{x^3}{3}-\cdots}.
\end{align*}
Esta construcción explica la forma de la serie; no es una fórmula aislada de memorización. En $x=1$ la serie converge por Leibniz. Además su suma es $\ln2$: la identidad geométrica finita produce
\[\frac1{1+t}=\sum_{n=0}^{N-1}(-t)^n+\frac{(-t)^N}{1+t},\qquad 0\le t\le1.\]
Al integrar entre $0$ y $1$, el resto tiene valor absoluto no mayor que $\int_0^1 t^N\,dt=1/(N+1)\to0$. Por ello $\sum_{n=1}^\infty(-1)^{n-1}/n=\ln2$. En $x=-1$ la serie diverge. La representación es válida en $(-1,1]$.
\end{ejemplo}
\begin{ejemplo}[De la geométrica a la arctangente]
Al sustituir $x$ por $-x^2$ en la geométrica,
\[\frac1{1+x^2}=\sum_{n=0}^\infty(-1)^n x^{2n},\qquad |x|<1.\]
Integrando desde $0$,
\[\arctan x=\boxed{\sum_{n=0}^\infty\frac{(-1)^n x^{2n+1}}{2n+1}}
=x-\frac{x^3}{3}+\frac{x^5}{5}-\cdots.\]
En $x=\pm1$ la serie converge condicionalmente por Leibniz y coincide con $\arctan(\pm1)=\pm\pi/4$.
\end{ejemplo}
\begin{ejemplo}[Derivar una serie de potencias]
Partimos de $\sum_{n=0}^\infty x^n=1/(1-x)$ para $|x|<1$. La derivada término a término ofrece
\[\sum_{n=1}^\infty nx^{n-1}=\frac1{(1-x)^2}.\]
Multiplicando por $x$,
\[\boxed{\sum_{n=1}^\infty nx^n=\frac{x}{(1-x)^2}},\qquad |x|<1.\]
En este intervalo está justificada la derivación; no es lícito extender la igualdad a $x=1$.
\end{ejemplo}

%=========================================================
\section{Polinomios de Taylor y de Maclaurin}
%=========================================================
La aproximación lineal $f(x)\approx f(a)+f'(a)(x-a)$ usa la posición y la pendiente de $f$ en $a$. Para lograr mayor precisión podemos exigir que un polinomio tenga, además, las mismas derivadas segunda, tercera, etc. que la función en el punto de referencia.

Sea
\[P_n(x)=c_0+c_1(x-a)+c_2(x-a)^2+\cdots+c_n(x-a)^n.\]
Si pedimos $P_n^{(k)}(a)=f^{(k)}(a)$ para $0\le k\le n$, obtenemos sucesivamente
\[c_0=f(a),\quad c_1=f'(a),\quad c_2=\frac{f''(a)}{2!},\quad\ldots,\quad c_n=\frac{f^{(n)}(a)}{n!}.\]
Esta construcción conduce a una definición natural.
\begin{definicion}[Polinomio de Taylor]
Si $f$ posee $n$ derivadas en $a$, su polinomio de Taylor de grado $n$ alrededor de $a$ es
\[\boxed{P_n(x)=\sum_{k=0}^n\frac{f^{(k)}(a)}{k!}(x-a)^k.}\]
Cuando $a=0$ se denomina \emph{polinomio de Maclaurin}.
\end{definicion}
\begin{ejemplo}[Maclaurin de $e^x$]
Como todas las derivadas de $f(x)=e^x$ son $e^x$ y $f^{(k)}(0)=1$,
\[P_0(x)=1,\qquad P_1(x)=1+x,\qquad
P_2(x)=1+x+\frac{x^2}{2},\]
\[\boxed{P_4(x)=1+x+\frac{x^2}{2}+\frac{x^3}{6}+\frac{x^4}{24}}.\]
La sucesión de polinomios aproxima cada vez mejor la función cerca de $0$; más adelante precisaremos hasta dónde puede garantizarse la convergencia.
\end{ejemplo}
\begin{center}
\begin{tikzpicture}
\begin{axis}[width=.78\linewidth,height=6.0cm,axis lines=middle,xmin=-1.4,xmax=1.4,ymin=-.1,ymax=4.25,domain=-1.35:1.35,samples=150,xlabel={$x$},ylabel={$y$},legend style={at={(.02,.98)},anchor=north west,font=\footnotesize,draw=none}]
\addplot[very thick,black] {exp(x)}; \addlegendentry{$e^x$}
\addplot[thick,red!80!black,dashed] {1+x+x^2/2};\addlegendentry{$P_2(x)$}
\addplot[thick,blue!80!black,dotted] {1+x+x^2/2+x^3/6+x^4/24};\addlegendentry{$P_4(x)$}
\end{axis}
\end{tikzpicture}

{\small Aproximaciones de Maclaurin de grados $2$ y $4$ a $e^x$ cerca de $x=0$.}
\end{center}

\begin{ejemplo}[Maclaurin de $\sin x$ y $\cos x$]
Las derivadas de $\sin x$ se repiten cada cuatro pasos: $\sin x,\cos x,-\sin x,-\cos x$. En $0$ las derivadas pares son cero y las impares alternan signos. Resulta
\[\boxed{\sin x=x-\frac{x^3}{3!}+\frac{x^5}{5!}-\frac{x^7}{7!}+\cdots}.\]
De forma análoga,
\[\boxed{\cos x=1-\frac{x^2}{2!}+\frac{x^4}{4!}-\frac{x^6}{6!}+\cdots}.\]
Estas series convergen para todo $x\in\R$, como se puede verificar mediante el criterio de la razón.
\end{ejemplo}

\subsection*{El término residual: una aproximación no es una identidad exacta}
Es crucial distinguir
\[f(x)=P_n(x)+R_n(x),\]
donde $R_n(x)$ representa el error de truncamiento. Conocer solamente el polinomio no determina el tamaño del error.
\begin{teorema}[Resto de Taylor en forma de Lagrange]
Supongamos que $f$ posee derivadas continuas hasta orden $n+1$ entre $a$ y $x$. Existe un punto $\xi$ estrictamente entre $a$ y $x$ tal que
\[\boxed{R_n(x)=\frac{f^{(n+1)}(\xi)}{(n+1)!}(x-a)^{n+1}.}\]
Si $|f^{(n+1)}(t)|\le M$ en todo el intervalo que une $a$ con $x$, entonces
\[\boxed{|R_n(x)|\le\frac{M}{(n+1)!}|x-a|^{n+1}.}\]
\end{teorema}
\begin{proof}[Justificación a partir del Teorema de Rolle]
Para $x\ne a$ escriba $f(x)-P_n(x)=K(x-a)^{n+1}$, donde $K$ es un número fijo determinado por $x$. Defina
\[G(t)=f(t)-P_n(t)-K(t-a)^{n+1}.\]
Por construcción, $G(a)=0$, $G(x)=0$ y, además, $G'(a)=\cdots=G^{(n)}(a)=0$, porque $P_n$ comparte las primeras $n$ derivadas de $f$ en $a$. Aplicando Rolle sucesivamente a $G$ y sus derivadas, existe un $\xi$ entre $a$ y $x$ tal que $G^{(n+1)}(\xi)=0$. Pero
\[G^{(n+1)}(\xi)=f^{(n+1)}(\xi)-(n+1)!K.\]
Por tanto, $K=f^{(n+1)}(\xi)/(n+1)!$. Sustituyendo en $f(x)-P_n(x)=K(x-a)^{n+1}$ queda la expresión del resto. La cota se obtiene tomando valores absolutos.
\end{proof}
\begin{ejemplo}[Aproximar $e^{1/2}$ y acotar el error]
Utilicemos el polinomio de Maclaurin de grado $3$:
\[P_3(1/2)=1+\frac12+\frac{(1/2)^2}{2}+\frac{(1/2)^3}{6}
=\boxed{\frac{79}{48}\approx1.6458333}.\]
Como la cuarta derivada de $e^x$ sigue siendo $e^x$, para $0\le t\le1/2$ se tiene $e^t\le e^{1/2}<2$. Por Lagrange,
\[|R_3(1/2)|\le\frac{2}{4!}\left(\frac12\right)^4
=\boxed{\frac1{192}\approx0.00521}.\]
Así tenemos una aproximación \emph{con garantía}: el verdadero valor difiere de $79/48$ en menos de $0.00521$.
\end{ejemplo}
\begin{ejemplo}[Aproximación de $\sin(0.2)$]
El polinomio $P_3(x)=x-x^3/6$ da
\[\sin(0.2)\approx0.2-\frac{0.2^3}{6}=\boxed{0.1986667}.\]
Si usamos el desarrollo hasta grado $4$ (cuyo coeficiente de $x^4$ es cero), el resto de Lagrange involucra la quinta derivada $\cos \xi$, de valor absoluto a lo más $1$. Por tanto,
\[|\sin(0.2)-P_4(0.2)|\le\frac{0.2^5}{5!}<2.7\times10^{-6}.\]
Esta cota muestra cuantitativamente por qué el procedimiento funciona cerca del centro.
\end{ejemplo}

%=========================================================
\section{Representaciones mediante series de Taylor y Maclaurin}
%=========================================================
Si el resto $R_n(x)$ tiende a cero cuando $n\to\infty$, la función coincide con la suma de su serie de Taylor:
\[f(x)=\sum_{n=0}^\infty\frac{f^{(n)}(a)}{n!}(x-a)^n.\]
Esta igualdad requiere justificar que el resto desaparece: la existencia de derivadas de todos los órdenes, por sí sola, no garantiza que la serie represente la función.
\begin{proposicion}
Las series de Maclaurin de $e^x$, $\sin x$ y $\cos x$ representan esas funciones para todo $x\in\R$.
\end{proposicion}
\begin{proof}
Fijado $x$, las derivadas de $e^t$ están acotadas por $e^{|x|}$ en el intervalo con extremos $0$ y $x$. Las de $\sin t$ y $\cos t$ están acotadas por $1$. Por la estimación de Lagrange, en los tres casos
\[|R_n(x)|\le M\frac{|x|^{n+1}}{(n+1)!},\]
con una constante $M$ independiente de $n$ (para ese $x$ fijo). La sucesión $|x|^{n+1}/(n+1)!$ tiende a cero por el criterio de la razón, de modo que $R_n(x)\to0$. Por ello la función coincide con su serie de Taylor.
\end{proof}

\begin{formula}[Series fundamentales]
\begin{align*}
\frac1{1-x}&=\sum_{n=0}^\infty x^n, &&|x|<1,\\
e^x&=\sum_{n=0}^\infty\frac{x^n}{n!}, &&x\in\R,\\
\sin x&=\sum_{n=0}^\infty\frac{(-1)^n x^{2n+1}}{(2n+1)!}, &&x\in\R,\\
\cos x&=\sum_{n=0}^\infty\frac{(-1)^n x^{2n}}{(2n)!}, &&x\in\R,\\
\ln(1+x)&=\sum_{n=1}^\infty\frac{(-1)^{n-1}x^n}{n}, &&-1<x\le1,\\
\arctan x&=\sum_{n=0}^\infty\frac{(-1)^n x^{2n+1}}{2n+1}, &&|x|\le1.
\end{align*}
En los extremos de las dos últimas series debe entenderse la convergencia puntual, examinada separadamente.
\end{formula}
\begin{ejemplo}[Centrar Taylor fuera del origen]
Encuentre el polinomio de Taylor de grado $2$ de $f(x)=\ln x$ alrededor de $a=1$. Dado que $f(1)=0$, $f'(x)=1/x$ y $f''(x)=-1/x^2$, se tiene
\[\boxed{P_2(x)=(x-1)-\frac12(x-1)^2}.\]
Por ejemplo, $\ln(1.1)\approx0.1-0.005=0.095$. Cambiar el centro puede mejorar considerablemente las aproximaciones cerca del valor de interés.
\end{ejemplo}
\begin{ejemplo}[Una función compuesta]
Para representar $e^{-x^2}$, sustituimos $t=-x^2$ en la serie de $e^t$:
\[e^{-x^2}=\sum_{n=0}^\infty\frac{(-1)^n x^{2n}}{n!}
=1-x^2+\frac{x^4}{2!}-\frac{x^6}{3!}+\cdots.\]
El criterio de la razón garantiza convergencia para todo $x\in\R$.
\end{ejemplo}

%=========================================================
\section{Aproximación de integrales mediante series}
%=========================================================
Hay integrales definidas cuyo integrando no tiene antiderivada elemental conocida, aunque puedan aproximarse con sumas, cuadraturas o series. El procedimiento de Taylor consiste en expandir el integrando, integrar un polinomio y acotar la contribución del resto.
\begin{ejemplo}[Una integral de tipo gaussiano]
Aproximemos
\[I=\int_0^{1/2}e^{-x^2}\dd.\]
De la serie anterior truncamos después del término de grado $6$:
\[e^{-x^2}\approx1-x^2+\frac{x^4}{2}-\frac{x^6}{6}.\]
Integramos término a término el polinomio:
\begin{align*}
I&\approx\left[x-\frac{x^3}{3}+\frac{x^5}{10}-\frac{x^7}{42}\right]_0^{1/2}\\
&=\frac12-\frac1{24}+\frac1{320}-\frac1{5376}
\approx\boxed{0.4612723}.
\end{align*}
En $0\le x\le1/2$, los términos de $e^{-x^2}$ son alternantes y sus magnitudes decrecen. El error puntual no supera el término siguiente $x^8/24$, por lo que el error de la integral está acotado por
\[|I-I_{\mathrm{aprox}}|\le\int_0^{1/2}\frac{x^8}{24}\dd
=\boxed{\frac{(1/2)^9}{216}<10^{-5}}.\]
Así no sólo encontramos una aproximación, sino una estimación cuantitativa de su precisión.
\end{ejemplo}
\begin{ejemplo}[Aproximar $\int_0^{0.2}\sin(x^2)\dd$]
Al sustituir $t=x^2$ en la serie de $\sin t$,
\[\sin(x^2)=x^2-\frac{x^6}{3!}+\frac{x^{10}}{5!}-\cdots.\]
Integrando los dos primeros términos,
\[I\approx\int_0^{0.2}\left(x^2-\frac{x^6}{6}\right)\dd
=\left[\frac{x^3}{3}-\frac{x^7}{42}\right]_0^{0.2}.\]
El resultado es aproximadamente $0.00266636$. El siguiente término integrado acota el error por
\[\int_0^{0.2}\frac{x^{10}}{120}\dd
=\frac{0.2^{11}}{1320}<2\times10^{-11}.\]
El propósito no es obtener una expresión elemental, sino aproximar de manera controlada.
\end{ejemplo}
\begin{observacion}
En un intervalo cerrado contenido estrictamente dentro del radio de convergencia de una serie de potencias, la convergencia es uniforme, lo cual justifica intercambiar integral y suma. Si el intervalo toca un extremo de convergencia, deben revisarse condiciones adicionales.
\end{observacion}

%=========================================================
\section{Visualización y verificación con R}
%=========================================================
Estas herramientas sirven para examinar tendencias y confirmar cálculos, no para reemplazar las pruebas de convergencia.
\begin{ejemplo}[Sucesión de sumas parciales]
Para la armónica y para la geométrica $\sum_{n=1}^\infty2^{-n}$:
\par\medskip
\begin{verbatim}
n <- 1:100
H <- cumsum(1/n)
G <- cumsum((1/2)^n)
par(mfrow = c(1, 2))
plot(n, H, type = "l", main = "Serie armonica",
     xlab = "N", ylab = "S_N")
plot(n, G, type = "l", main = "Serie geometrica",
     xlab = "N", ylab = "S_N")
abline(h = 1, lty = 2)
par(mfrow = c(1, 1))
\end{verbatim}
La primera suma crece sin cota (lentamente), mientras que la segunda se aproxima a $1$. Una gráfica no demuestra la convergencia, pero hace visible la diferencia.
\end{ejemplo}
\begin{ejemplo}[Comparación entre $e^x$ y sus polinomios]
El siguiente código grafica la función exponencial y dos polinomios:

\par\medskip
\begin{verbatim}
x <- seq(-1.5,1.5,.01)
P2 <- 1 + x + x^2/2
P4 <- P2 + x^3/6 + x^4/24
plot(x, exp(x), type = "l", lwd = 2,
     xlab = "x", ylab = "y", main = "Maclaurin de exp(x)")
lines(x, P2, lty = 2)
lines(x, P4, lty = 3)
legend("topleft", legend = c("exp(x)", "P2", "P4"),
       lty = c(1, 2, 3), bty = "n")
\end{verbatim}
La coincidencia es más cercana alrededor de $x=0$, que es el centro de los desarrollos.
\end{ejemplo}
\Needspace{15\baselineskip}
\begin{ejemplo}[Verificar una integral aproximada]
La cuadratura numérica proporciona una verificación independiente:

\par\medskip
\begin{verbatim}
I_aprox <- 1/2 - 1/24
I_aprox <- I_aprox + 1/320 - 1/5376
I_num <- integrate(function(x) exp(-x^2), 0, 0.5)$value
cat("Taylor:", I_aprox, "\n")
cat("R:", I_num, "\n")
cat("Error:", abs(I_num - I_aprox), "\n")
\end{verbatim}
El error numérico debe satisfacer la cota del ejemplo anterior. Conviene confrontar siempre un cálculo numérico con una estimación teórica de error.
\end{ejemplo}

%=========================================================
\section{Banco de ejercicios}
%=========================================================
Resuelva primero sin software. Para cada serie, elija y justifique un criterio: no basta escribir ``converge'' o ``diverge''. Si aplica razón o raíz y obtiene límite $1$, cambie de herramienta. Si se pide intervalo de convergencia, examine ambos extremos por separado. Los ejercicios de Taylor deben acompañarse de una estimación del error cuando así se indique.

\begin{Ejer}\textbf{Sucesiones: cálculo y demostración de límites.}
\begin{enumerate}[series=banco4]
\item Determine $\lim_{n\to\infty}(4n+1)/(2n-3)$.
\item Determine $\lim_{n\to\infty}(n^2+3n)/(3n^2+1)$.
\item Determine $\lim_{n\to\infty}\sqrt{n^2+1}/n$.
\item Determine $\lim_{n\to\infty}(-1)^n/n$.
\item Analice la convergencia de $a_n=(-1)^n$.
\item Demuestre desde la definición que $a_n=3+2/n$ tiende a $3$.
\item Determine $\lim_{n\to\infty}n/(n+1)$ y proponga una elección de $N$ dado $\varepsilon>0$.
\item Determine $\lim_{n\to\infty}(2n^3-1)/(n^3+n)$.
\item Sea $a_1=1$ y $a_{n+1}=(a_n+2)/2$. Demuestre que $a_n$ es creciente, está acotada y encuentre su límite.
\item Explique por qué una sucesión acotada no necesariamente converge, y dé un ejemplo concreto.
\end{enumerate}
\end{Ejer}

\begin{Ejer}\textbf{Series geométricas y telescópicas.} Calcule la suma cuando exista.
\begin{enumerate}[resume=banco4]
\item $\displaystyle\sum_{n=0}^\infty (1/2)^n$.
\item $\displaystyle\sum_{n=0}^\infty 3(1/4)^n$.
\item $\displaystyle\sum_{n=1}^\infty 2/3^n$.
\item $\displaystyle\sum_{n=0}^\infty (-1/2)^n$.
\item $\displaystyle\sum_{n=1}^\infty (-1)^{n-1}/2^n$.
\item $\displaystyle\sum_{n=1}^\infty 1/[n(n+1)]$.
\item $\displaystyle\sum_{n=1}^\infty 1/[(n+1)(n+2)]$.
\item $\displaystyle\sum_{n=1}^\infty [1/n-1/(n+2)]$.
\item $\displaystyle\sum_{n=0}^\infty 2^n/3^{n+1}$.
\item $\displaystyle\sum_{n=1}^\infty (4/3)^n$: indique por qué no converge.
\end{enumerate}
\end{Ejer}

\begin{Ejer}\textbf{Prueba del término general y series $p$.} Decida convergencia o divergencia.
\begin{enumerate}[resume=banco4]
\item $\displaystyle\sum_{n=1}^\infty 1/n^2$.
\item $\displaystyle\sum_{n=1}^\infty 1/n$.
\item $\displaystyle\sum_{n=1}^\infty 1/\sqrt n$.
\item $\displaystyle\sum_{n=1}^\infty 1/n^{3/2}$.
\item $\displaystyle\sum_{n=1}^\infty n/(n+1)$.
\item $\displaystyle\sum_{n=1}^\infty (2n+1)/(3n-1)$.
\item $\displaystyle\sum_{n=1}^\infty (-1)^n$.
\item $\displaystyle\sum_{n=1}^\infty n^{-5/4}$.
\item $\displaystyle\sum_{n=1}^\infty (n+2)^{-1}$.
\item Para qué valores de $p$ converge $\displaystyle\sum_{n=1}^\infty 1/n^p$? Explique su relación con una integral impropia.
\end{enumerate}
\end{Ejer}

\begin{Ejer}\textbf{Comparación directa y por el límite.} Declare una serie de referencia.
\begin{enumerate}[resume=banco4]
\item $\displaystyle\sum_{n=1}^\infty 1/(n^2+4)$.
\item $\displaystyle\sum_{n=1}^\infty 1/(n+3)$.
\item $\displaystyle\sum_{n=1}^\infty (2n+1)/(n^3+1)$.
\item $\displaystyle\sum_{n=1}^\infty (n+2)/(n^2+1)$.
\item $\displaystyle\sum_{n=1}^\infty 1/\sqrt{n^2+1}$.
\item $\displaystyle\sum_{n=1}^\infty 1/\sqrt{n^3+1}$.
\item $\displaystyle\sum_{n=1}^\infty (3n^2+1)/(n^4+2)$.
\item $\displaystyle\sum_{n=1}^\infty (n^2+1)/(n^3+3)$.
\item $\displaystyle\sum_{n=1}^\infty 1/(n^2+\sqrt n)$.
\item Justifique que la serie $\sum_{n=1}^\infty (n^2+n)/(n^4+1)$ converge.
\end{enumerate}
\end{Ejer}

\begin{Ejer}\textbf{Series alternantes y convergencia absoluta.}
\begin{enumerate}[resume=banco4]
\item $\displaystyle\sum_{n=1}^\infty (-1)^{n-1}/n$.
\item $\displaystyle\sum_{n=1}^\infty (-1)^{n-1}/n^2$.
\item $\displaystyle\sum_{n=1}^\infty (-1)^n/\sqrt n$.
\item $\displaystyle\sum_{n=1}^\infty (-1)^n n/(n+1)$.
\item $\displaystyle\sum_{n=1}^\infty (-1)^{n-1}/(2n+1)$.
\item Para $\sum_{n=1}^\infty(-1)^{n-1}/n$, ¿cuántos términos garantizan un error no mayor que $10^{-2}$?
\item Determine si $\sum_{n=1}^\infty(-1)^n/(n+2)^{3/2}$ converge absolutamente o condicionalmente.
\item Explique por qué la serie alternante $\sum_{n=1}^\infty(-1)^n$ no cumple Leibniz.
\end{enumerate}
\end{Ejer}

\begin{Ejer}\textbf{Criterio de la razón.} Calcule el límite y explique si el criterio decide.
\begin{enumerate}[resume=banco4]
\item $\displaystyle\sum_{n=1}^\infty n/2^n$.
\item $\displaystyle\sum_{n=1}^\infty n!/4^n$.
\item $\displaystyle\sum_{n=1}^\infty 3^n/n!$.
\item $\displaystyle\sum_{n=1}^\infty n^2/5^n$.
\item $\displaystyle\sum_{n=1}^\infty (n!)^2/(2n)!$.
\item $\displaystyle\sum_{n=1}^\infty n!/n^n$.
\item $\displaystyle\sum_{n=1}^\infty 1/n^2$ (explique por qué $L=1$ es inconcluso).
\item $\displaystyle\sum_{n=1}^\infty 2^n/n^3$.
\end{enumerate}
\end{Ejer}

\begin{Ejer}\textbf{Criterio de la raíz.} Aplíquelo a potencias completas del término general.
\begin{enumerate}[resume=banco4]
\item $\displaystyle\sum_{n=1}^\infty (2/3)^n$.
\item $\displaystyle\sum_{n=1}^\infty [n/(2n+1)]^n$.
\item $\displaystyle\sum_{n=1}^\infty [(3n+1)/(2n+1)]^n$.
\item $\displaystyle\sum_{n=1}^\infty [(n+1)/(2n)]^n$.
\item $\displaystyle\sum_{n=1}^\infty [(n+1)/n]^n$.
\item $\displaystyle\sum_{n=1}^\infty [1/(n+1)]^n$.
\item Explique por qué la raíz da $L=1$ para $\sum1/n^2$ y qué otro criterio es apropiado.
\item Compare raíz y razón para $\sum_{n=1}^\infty(1/4)^n$.
\end{enumerate}
\end{Ejer}

\begin{Ejer}\textbf{Radio e intervalo de convergencia.} Analice siempre los dos extremos.
\begin{enumerate}[resume=banco4]
\item $\displaystyle\sum_{n=0}^\infty x^n$.
\item $\displaystyle\sum_{n=0}^\infty (x-2)^n/3^n$.
\item $\displaystyle\sum_{n=1}^\infty x^n/n$.
\item $\displaystyle\sum_{n=1}^\infty (-1)^{n-1}x^n/n$.
\item $\displaystyle\sum_{n=1}^\infty (x-1)^n/(n2^n)$.
\item $\displaystyle\sum_{n=1}^\infty (x+2)^n/(n3^n)$.
\item $\displaystyle\sum_{n=0}^\infty n!x^n$.
\item $\displaystyle\sum_{n=0}^\infty x^n/n!$.
\item $\displaystyle\sum_{n=1}^\infty x^n/n^2$.
\item $\displaystyle\sum_{n=1}^\infty n(x-2)^n/4^n$.
\end{enumerate}
\end{Ejer}

\begin{Ejer}\textbf{Derivar e integrar series de potencias.} Indique el intervalo en que la operación es válida.
\begin{enumerate}[resume=banco4]
\item Derive $\sum_{n=0}^\infty x^n$ para $|x|<1$ y obtenga una fórmula para $\sum_{n=1}^\infty nx^{n-1}$.
\item Utilice el resultado anterior para calcular $\sum_{n=1}^\infty nx^n$ cuando $|x|<1$.
\item Integre $\sum_{n=0}^\infty x^n$ entre $0$ y $x$ para deducir una serie de $-\ln(1-x)$.
\item Integre la serie de $1/(1+x)$ entre $0$ y $x$ para obtener $\ln(1+x)$.
\item Encuentre la serie de potencias de $1/(1+x^2)$ alrededor de $0$.
\item Integre la serie anterior para obtener la representación de $\arctan x$.
\item Derive término a término $\sum_{n=1}^\infty x^n/n$ dentro de su intervalo de convergencia.
\item Explique por qué no es válido derivar la serie geométrica en $x=1$.
\end{enumerate}
\end{Ejer}

\begin{Ejer}\textbf{Polinomios de Taylor y Maclaurin.} Muestre las derivadas utilizadas.
\begin{enumerate}[resume=banco4]
\item Obtenga el polinomio de Maclaurin de grado $3$ para $e^x$.
\item Obtenga el polinomio de Maclaurin de grado $5$ para $\sin x$.
\item Obtenga el polinomio de Maclaurin de grado $4$ para $\cos x$.
\item Obtenga el polinomio de Maclaurin de grado $3$ para $\ln(1+x)$.
\item Obtenga el polinomio de Taylor de grado $2$ para $\ln x$ alrededor de $a=1$.
\item Obtenga el polinomio de Taylor de grado $3$ para $e^x$ alrededor de $a=1$.
\item Obtenga el polinomio de Taylor de grado $2$ de $f(x)=1/x$ alrededor de $a=1$.
\item Obtenga el polinomio de Maclaurin de grado $4$ para $e^{-x^2}$.
\item Obtenga el polinomio de Maclaurin de grado $6$ para $\cos(2x)$.
\item Obtenga el polinomio de Maclaurin de grado $5$ para $\sin(x^2)$ (aclare qué grados aparecen realmente).
\end{enumerate}
\end{Ejer}

\begin{Ejer}\textbf{Aproximaciones y cotas de error.} Identifique la derivada que controla el resto.
\begin{enumerate}[resume=banco4]
\item Aproximar $e^{0.1}$ con Maclaurin de grado $2$ y acotar el error por Lagrange.
\item Aproximar $e^{0.5}$ con Maclaurin de grado $3$ y justificar una cota numérica.
\item Aproximar $\sin(0.2)$ con un polinomio de grado $3$ y controlar el resto.
\item Aproximar $\cos(0.1)$ con un polinomio de grado $4$ y controlar el resto.
\item Aproximar $\ln(1.1)$ mediante Taylor de grado $2$ centrado en $1$.
\item Aproximar $\sqrt{1.1}$ mediante Taylor de grado $2$ de $\sqrt{1+x}$ centrado en $x=0$.
\item ¿Cuántos términos de la serie de $e^x$ bastan para aproximar $e^{0.2}$ con un error menor que $10^{-4}$? Justifique.
\item Explique por qué el tamaño de $|x-a|$ influye en la cota de Lagrange.
\end{enumerate}
\end{Ejer}

\begin{Ejer}\textbf{Integrales mediante series.} Desarrolle el integrando, integre y estime el error.
\begin{enumerate}[resume=banco4]
\item Aproximar $\displaystyle\int_0^{0.5}e^{-x^2}\dd$ usando cuatro términos de Maclaurin.
\item Aproximar $\displaystyle\int_0^{0.2}\sin(x^2)\dd$ usando dos términos de su serie.
\item Aproximar $\displaystyle\int_0^{0.5}\cos(x^2)\dd$ usando tres términos de su serie.
\item Aproximar $\displaystyle\int_0^{0.3}\frac1{1+x^2}\dd$ integrando la serie geométrica apropiada.
\item Usar la serie de $e^x$ para aproximar $\displaystyle\int_0^{0.2}e^x\dd$ con tres términos.
\item Usar la serie de $\sin x$ para aproximar $\displaystyle\int_0^{0.4}\sin x\dd$ con tres términos.
\item Deducir una serie de potencias para $\displaystyle\int_0^x\frac1{1+t^2}\,dt$ y reconocer la función obtenida.
\item Justificar por qué integrar término a término una serie de potencias es válido en un intervalo cerrado situado dentro de su radio de convergencia.
\end{enumerate}
\end{Ejer}

\section*{Formulario breve y guía para seleccionar el criterio}
\begin{formula}[Reglas de decisión]
\begin{tabular}{@{}p{.37\linewidth}p{.53\linewidth}@{}}
\toprule
\textbf{Estructura} & \textbf{Prueba inicial sugerida}\\
\midrule
$\sum ar^n$ & Geométrica: $|r|<1$\\
$\sum n^{-p}$ & Serie $p$: $p>1$\\
$a_n\not\to0$ & Divergencia inmediata\\
Términos positivos semejantes a $n^{-p}$ & Comparación directa o por límite\\
Signos alternantes & Leibniz; después estudiar $\sum|a_n|$\\
Factoriales o exponenciales & Razón\\
Todo el término elevado a $n$ & Raíz\\
Potencias de $(x-c)$ & Razón/raíz para $R$ y prueba independiente de extremos\\
Aproximación local de $f(x)$ & Taylor con cota de $R_n$\\
\bottomrule
\end{tabular}
\end{formula}

\end{document}
